\chapter{Requirements Analysis and Specification}

\section*{Introduction}

This chapter formalises the system requirements as functional and non-functional
needs, identifies the actors, and presents the UML models that guide the design
phase. The analysis drew on the comparative study from the previous chapter,
supplemented by concrete usage scenarios.

\section{Actor Identification}

\begin{description}[leftmargin=2.5cm, style=nextline]
  \item[\textbf{Administrator}] Most privileged role. Manages organisations, users,
    servers and platform configuration. Accesses the Admin API (port 8001).
  \item[\textbf{Operator}] Technical user who consults dashboards, manages runbooks,
    configures alert rules and interprets AI analyses. Role \texttt{devops} in the database.
  \item[\textbf{Viewer}] Read-only access. Consults dashboards and PDF reports. Designed
    for non-technical stakeholders.
  \item[\textbf{System Agent}] Non-human actor. The \texttt{caimp-agent} daemon process
    runs on monitored servers, authenticates via mTLS certificate and sends metrics via OTLP.
\end{description}

\section{Functional Requirements}

\begin{enumerate}[label=\textbf{FR\arabic*}, leftmargin=2.2cm, itemsep=4pt]
  \item The agent must collect CPU, memory and disk metrics every 15 seconds and transmit them via OTLP/HTTP.
  \item The agent must self-register with the platform using a single-use enrollment token.
  \item The agent must use an mTLS certificate issued by the internal CA for every connection.
  \item The agent must detect and report Docker events automatically.
  \item The agent must monitor package installation logs and SSH events from authentication logs.
  \item The platform must detect anomalies in real time using at least three distinct algorithms.
  \item Detected anomalies must be published on NATS JetStream for asynchronous processing.
  \item The AI analysis engine must query the vector database to retrieve similar past incidents and relevant runbooks.
  \item The AI analysis engine must submit a structured prompt to Ollama and validate the JSON response.
  \item The AI explanation must contain: natural-language explanation, probable root cause, recommended action, and confidence level.
  \item The correlation engine must query \texttt{change\_events} within the 30-minute window before each anomaly.
  \item The correlation engine must score each change by temporal proximity and type, and produce a ranked list of probable causes.
  \item The forecast engine must compute Holt-Winters predictions over 24 hours for each server/metric pair.
  \item The forecast engine must calculate the \textit{time to critical} — the predicted delay before threshold breach.
  \item Alert deduplication must suppress identical alerts within 15 minutes.
  \item The query API must expose a \texttt{/dashboard/summary} endpoint aggregating global status, potential issues and recommended actions.
  \item The user interface must present the global status, potential issues sorted by criticality and numbered recommended actions on the home page.
  \item The platform must allow runbook management (create, edit, delete) with automatic embedding re-indexing.
  \item The report module must generate a multi-section PDF report readable by non-technical stakeholders.
  \item The interface must provide an SSE chatbot for natural-language questions about incidents and servers.
  \item The platform must receive GitHub and GitLab webhooks for deployment events.
  \item The administrator must be able to configure notification channels (Slack, PagerDuty).
  \item Status updates must be pushed in real time to connected browsers via WebSocket.
  \item The platform must support multi-tenant isolation: one organisation's data must never be accessible by another.
  \item Each service must provide a health endpoint used by Docker healthchecks.
\end{enumerate}

\section{Non-Functional Requirements}

\textbf{Performance:} OTLP batch processing $<$50\,ms latency p50; throughput
$\geq$1,000 metrics/second; API GET responses $<$300\,ms for 50 servers; full AI
explanation $<$30 seconds.

\textbf{Security:} Agent mTLS (TLS 1.3); single-use enrollment tokens expiring
in 30 minutes; JWT HS256 with 15-minute expiry; RLS PostgreSQL; AES-256-GCM
encrypted Splunk secrets; immutable audit log.

\textbf{Availability:} 99.5\% uptime target; Docker healthchecks with
\texttt{unless-stopped} restart; local SQLite buffering in agent; JetStream
message durability on disk.

\textbf{Maintainability:} Independent, individually deployable services;
versioned DB migrations; conventional commits; inline documentation.

\section{Use Case Diagram}

\begin{figure}[H]
\centering
\begin{tikzpicture}[
  actor/.style={draw=NavyBlue, fill=NavyBlue!10, rounded corners=3pt,
    minimum width=1.8cm, minimum height=0.7cm, font=\small\bfseries},
  uc/.style={draw=TealBlue, fill=TealBlue!8, ellipse,
    minimum width=3.4cm, minimum height=0.9cm, font=\small, align=center},
  sysbox/.style={draw=NavyBlue!50, dashed, rounded corners=6pt, fill=NavyBlue!3}
]
\node[sysbox, minimum width=13cm, minimum height=11cm] (sys) at (5.5,0) {};
\node[above, font=\small\bfseries\color{NavyBlue}] at (sys.north) {CAIMP System};

\node[actor] (admin)  at (-1.2, 3.0) {Admin};
\node[actor] (ops)    at (-1.2, 0.5) {Operator};
\node[actor] (viewer) at (-1.2,-2.0) {Viewer};
\node[actor] (agent)  at (12.5, 0.5) {Agent};

\node[uc] (login)    at (3.0, 4.5) {Authenticate};
\node[uc] (usrmgr)   at (7.0, 4.5) {Manage users};
\node[uc] (srvmgr)   at (7.0, 3.0) {Manage servers};
\node[uc] (dash)     at (3.0, 1.8) {View dashboard};
\node[uc] (alerts)   at (7.0, 1.8) {View live alerts};
\node[uc] (runbooks) at (3.0, 0.2) {Manage runbooks};
\node[uc] (chat)     at (7.0, 0.2) {Use AI chat};
\node[uc] (report)   at (3.0,-1.4) {Export PDF};
\node[uc] (forecast) at (7.0,-1.4) {View forecasts};
\node[uc] (metrics)  at (5.0,-3.0) {Send metrics\\(OTLP)};
\node[uc] (enroll)   at (9.0,-3.0) {Enrol (mTLS)};

\draw[->, NavyBlue] (admin)  -- (login);
\draw[->, NavyBlue] (admin)  -- (usrmgr);
\draw[->, NavyBlue] (admin)  -- (srvmgr);
\draw[->, NavyBlue] (ops)    -- (login);
\draw[->, NavyBlue] (ops)    -- (dash);
\draw[->, NavyBlue] (ops)    -- (alerts);
\draw[->, NavyBlue] (ops)    -- (runbooks);
\draw[->, NavyBlue] (ops)    -- (chat);
\draw[->, NavyBlue] (ops)    -- (forecast);
\draw[->, NavyBlue] (viewer) -- (login);
\draw[->, NavyBlue] (viewer) -- (dash);
\draw[->, NavyBlue] (viewer) -- (report);
\draw[->, NavyBlue] (agent)  -- (metrics);
\draw[->, NavyBlue] (agent)  -- (enroll);
\end{tikzpicture}
\caption{Global use case diagram of the \caimp{} platform}
\label{fig:ucd}
\end{figure}

\section{Requirement Traceability Matrix}

\begin{table}[H]
\centering
\caption{Traceability matrix: functional requirements — components}
\label{tab:traceability}
\small
\begin{tabularx}{\textwidth}{lXll}
\toprule
\textbf{FR} & \textbf{Requirement} & \textbf{Component} & \textbf{Module} \\
\midrule
FR1  & OTLP metric collection       & \servicename{caimp-agent}        & \texttt{collectors.py} \\
FR6  & 3-algorithm anomaly detection & \servicename{telemetry-writer}   & \texttt{detector.go} \\
FR8--FR10 & AI explanation           & \servicename{ai-worker}          & \texttt{prompt\_builder.py} \\
FR11--FR12 & Change correlation      & \servicename{correlation-engine} & \texttt{main.py} \\
FR13--FR14 & Holt-Winters forecast   & \servicename{forecast-engine}    & \texttt{forecaster.py} \\
FR15 & Alert deduplication           & \servicename{alert-engine}       & \texttt{worker.py} \\
FR16--FR17 & Answers-first dashboard & \servicename{query-api}          & \texttt{dashboard.py} \\
FR18 & Runbook CRUD + RAG            & \servicename{query-api}          & \texttt{incidents.py} \\
FR21 & GitHub/GitLab webhooks        & \servicename{change-tracker}     & \texttt{main.py} \\
FR23 & WebSocket push                & \servicename{ws-gateway}         & \texttt{main.py} \\
FR24 & Multi-tenant isolation        & PostgreSQL                       & RLS migrations \\
\bottomrule
\end{tabularx}
\end{table}

\section*{Conclusion}

The requirements analysis identified 25 functional requirements covering the full
scope of \caimp{}, and formalised non-functional requirements around performance,
security, availability and maintainability. The UML models presented constitute
the contract between analysis and the design covered in the next chapter.
